\section*{原书目录（原书页码）}\phantomsection\label{front:source-contents}
本节及后面的原书表、图、例目录保留原书页码；点击条目可跳到本译稿的对应位置. 译稿自己的页码见后面的“中文译稿目录”.

\begingroup\small
\begin{longtable}{@{}p{.86\textwidth}r@{}}
\hyperref[front:preface-second]{第二版前言} & \hyperref[front:preface-second]{xi} \\
\hyperref[front:preface-first]{第一版前言} & \hyperref[front:preface-first]{xiii} \\
\hyperref[front:authors]{作者} & \hyperref[front:authors]{xvii} \\
\hyperref[front:tables]{表目录} & \hyperref[front:tables]{xix} \\
\hyperref[front:figures]{图目录} & \hyperref[front:figures]{xxi} \\
\hyperref[front:examples]{例目录} & \hyperref[front:examples]{xxiii} \\
\hyperref[sec:1]{1 概率论} & \hyperref[sec:1]{1} \\
\hyperref[sec:1.1]{1.1 集合论} & \hyperref[sec:1.1]{1} \\
\hyperref[sec:1.2]{1.2 概率论} & \hyperref[sec:1.2]{4} \\
\hyperref[sec:1.2.1]{1.2.1 公理基础} & \hyperref[sec:1.2.1]{5} \\
\hyperref[sec:1.2.2]{1.2.2 概率运算} & \hyperref[sec:1.2.2]{8} \\
\hyperref[sec:1.2.3]{1.2.3 计数} & \hyperref[sec:1.2.3]{11} \\
\hyperref[sec:1.2.4]{1.2.4 枚举结果} & \hyperref[sec:1.2.4]{13} \\
\hyperref[sec:1.3]{1.3 条件概率与独立性} & \hyperref[sec:1.3]{16} \\
\hyperref[sec:1.4]{1.4 随机变量} & \hyperref[sec:1.4]{22} \\
\hyperref[sec:1.5]{1.5 分布函数} & \hyperref[sec:1.5]{24} \\
\hyperref[sec:1.6]{1.6 密度函数与质量函数} & \hyperref[sec:1.6]{28} \\
\hyperref[sec:1.7]{1.7 习题} & \hyperref[sec:1.7]{31} \\
\hyperref[sec:1.8]{1.8 杂录} & \hyperref[sec:1.8]{37} \\
\hyperref[sec:2]{2 变换与期望} & \hyperref[sec:2]{38} \\
\hyperref[sec:2.1]{2.1 随机变量的函数的分布} & \hyperref[sec:2.1]{38} \\
\hyperref[sec:2.2]{2.2 期望值} & \hyperref[sec:2.2]{45} \\
\hyperref[sec:2.3]{2.3 矩与矩母函数} & \hyperref[sec:2.3]{48} \\
\hyperref[sec:2.4]{2.4 积分号下求导} & \hyperref[sec:2.4]{56} \\
\hyperref[sec:2.5]{2.5 习题} & \hyperref[sec:2.5]{62} \\
\hyperref[sec:2.6]{2.6 杂录} & \hyperref[sec:2.6]{68} \\
\hyperref[sec:3]{3 常见分布族} & \hyperref[sec:3]{70} \\
\hyperref[sec:3.1]{3.1 引言} & \hyperref[sec:3.1]{70} \\
\hyperref[sec:3.2]{3.2 离散分布} & \hyperref[sec:3.2]{70} \\
\hyperref[sec:3.3]{3.3 连续分布} & \hyperref[sec:3.3]{81} \\
\hyperref[sec:3.4]{3.4 指数族} & \hyperref[sec:3.4]{90} \\
\hyperref[sec:3.5]{3.5 位置族与尺度族} & \hyperref[sec:3.5]{95} \\
\hyperref[sec:3.6]{3.6 不等式与恒等式} & \hyperref[sec:3.6]{100} \\
\hyperref[sec:3.6.1]{3.6.1 概率不等式} & \hyperref[sec:3.6.1]{100} \\
\hyperref[sec:3.6.2]{3.6.2 恒等式} & \hyperref[sec:3.6.2]{102} \\
\hyperref[sec:3.7]{3.7 习题} & \hyperref[sec:3.7]{105} \\
\hyperref[sec:3.8]{3.8 杂录} & \hyperref[sec:3.8]{111} \\
\hyperref[sec:4]{4 多个随机变量} & \hyperref[sec:4]{114} \\
\hyperref[sec:4.1]{4.1 联合分布与边缘分布} & \hyperref[sec:4.1]{114} \\
\hyperref[sec:4.2]{4.2 条件分布与独立性} & \hyperref[sec:4.2]{121} \\
\hyperref[sec:4.3]{4.3 二元变换} & \hyperref[sec:4.3]{128} \\
\hyperref[sec:4.4]{4.4 分层模型与混合分布} & \hyperref[sec:4.4]{133} \\
\hyperref[sec:4.5]{4.5 协方差与相关} & \hyperref[sec:4.5]{138} \\
\hyperref[sec:4.6]{4.6 多元分布} & \hyperref[sec:4.6]{145} \\
\hyperref[sec:4.7]{4.7 不等式} & \hyperref[sec:4.7]{152} \\
\hyperref[sec:4.7.1]{4.7.1 数值不等式} & \hyperref[sec:4.7.1]{152} \\
\hyperref[sec:4.7.2]{4.7.2 函数不等式} & \hyperref[sec:4.7.2]{155} \\
\hyperref[sec:4.8]{4.8 习题} & \hyperref[sec:4.8]{157} \\
\hyperref[sec:4.9]{4.9 杂录} & \hyperref[sec:4.9]{166} \\
\hyperref[sec:5]{5 随机样本的性质} & \hyperref[sec:5]{169} \\
\hyperref[sec:5.1]{5.1 随机样本的基本概念} & \hyperref[sec:5.1]{169} \\
\hyperref[sec:5.2]{5.2 随机样本中的随机变量之和} & \hyperref[sec:5.2]{172} \\
\hyperref[sec:5.3]{5.3 从正态分布抽样} & \hyperref[sec:5.3]{177} \\
\hyperref[sec:5.3.1]{5.3.1 样本均值与方差的性质} & \hyperref[sec:5.3.1]{178} \\
\hyperref[sec:5.3.2]{5.3.2 导出分布：Student的$t$与Snedecor的$F$} & \hyperref[sec:5.3.2]{181} \\
\hyperref[sec:5.4]{5.4 次序统计量} & \hyperref[sec:5.4]{184} \\
\hyperref[sec:5.5]{5.5 收敛概念} & \hyperref[sec:5.5]{188} \\
\hyperref[sec:5.5.1]{5.5.1 依概率收敛} & \hyperref[sec:5.5.1]{189} \\
\hyperref[sec:5.5.2]{5.5.2 几乎处处收敛} & \hyperref[sec:5.5.2]{190} \\
\hyperref[sec:5.5.3]{5.5.3 依分布收敛} & \hyperref[sec:5.5.3]{191} \\
\hyperref[sec:5.5.4]{5.5.4 Delta方法} & \hyperref[sec:5.5.4]{195} \\
\hyperref[sec:5.6]{5.6 随机样本的生成} & \hyperref[sec:5.6]{200} \\
\hyperref[sec:5.6.1]{5.6.1 直接方法} & \hyperref[sec:5.6.1]{201} \\
\hyperref[sec:5.6.2]{5.6.2 间接方法} & \hyperref[sec:5.6.2]{204} \\
\hyperref[sec:5.6.3]{5.6.3 接受／拒绝算法} & \hyperref[sec:5.6.3]{206} \\
\hyperref[sec:5.7]{5.7 习题} & \hyperref[sec:5.7]{208} \\
\hyperref[sec:5.8]{5.8 杂录} & \hyperref[sec:5.8]{217} \\
\hyperref[sec:6]{6 数据约简原则} & \hyperref[sec:6]{221} \\
\hyperref[sec:6.1]{6.1 引言} & \hyperref[sec:6.1]{221} \\
\hyperref[sec:6.2]{6.2 充分性原则} & \hyperref[sec:6.2]{221} \\
\hyperref[sec:6.2.1]{6.2.1 充分统计量} & \hyperref[sec:6.2.1]{222} \\
\hyperref[sec:6.2.2]{6.2.2 最小充分统计量} & \hyperref[sec:6.2.2]{228} \\
\hyperref[sec:6.2.3]{6.2.3 辅助统计量} & \hyperref[sec:6.2.3]{230} \\
\hyperref[sec:6.2.4]{6.2.4 充分、辅助与完备统计量} & \hyperref[sec:6.2.4]{231} \\
\hyperref[sec:6.3]{6.3 似然原则} & \hyperref[sec:6.3]{236} \\
\hyperref[sec:6.3.1]{6.3.1 似然函数} & \hyperref[sec:6.3.1]{236} \\
\hyperref[sec:6.3.2]{6.3.2 形式化的似然原则} & \hyperref[sec:6.3.2]{238} \\
\hyperref[sec:6.4]{6.4 等变性原则} & \hyperref[sec:6.4]{241} \\
\hyperref[sec:6.5]{6.5 习题} & \hyperref[sec:6.5]{244} \\
\hyperref[sec:6.6]{6.6 杂录} & \hyperref[sec:6.6]{250} \\
\hyperref[sec:7]{7 点估计} & \hyperref[sec:7]{252} \\
\hyperref[sec:7.1]{7.1 引言} & \hyperref[sec:7.1]{252} \\
\hyperref[sec:7.2]{7.2 寻找估计量的方法} & \hyperref[sec:7.2]{253} \\
\hyperref[sec:7.2.1]{7.2.1 矩估计法} & \hyperref[sec:7.2.1]{253} \\
\hyperref[sec:7.2.2]{7.2.2 最大似然估计量} & \hyperref[sec:7.2.2]{255} \\
\hyperref[sec:7.2.3]{7.2.3 贝叶斯估计量} & \hyperref[sec:7.2.3]{262} \\
\hyperref[sec:7.2.4]{7.2.4 EM算法} & \hyperref[sec:7.2.4]{264} \\
\hyperref[sec:7.3]{7.3 评价估计量的方法} & \hyperref[sec:7.3]{267} \\
\hyperref[sec:7.3.1]{7.3.1 均方误差} & \hyperref[sec:7.3.1]{267} \\
\hyperref[sec:7.3.2]{7.3.2 最优无偏估计量} & \hyperref[sec:7.3.2]{270} \\
\hyperref[sec:7.3.3]{7.3.3 充分性与无偏性} & \hyperref[sec:7.3.3]{277} \\
\hyperref[sec:7.3.4]{7.3.4 损失函数意义下的最优性} & \hyperref[sec:7.3.4]{282} \\
\hyperref[sec:7.4]{7.4 习题} & \hyperref[sec:7.4]{287} \\
\hyperref[sec:7.5]{7.5 杂录} & \hyperref[sec:7.5]{298} \\
\hyperref[sec:8]{8 假设检验} & \hyperref[sec:8]{302} \\
\hyperref[sec:8.1]{8.1 引言} & \hyperref[sec:8.1]{302} \\
\hyperref[sec:8.2]{8.2 寻找检验的方法} & \hyperref[sec:8.2]{303} \\
\hyperref[sec:8.2.1]{8.2.1 似然比检验} & \hyperref[sec:8.2.1]{303} \\
\hyperref[sec:8.2.2]{8.2.2 贝叶斯检验} & \hyperref[sec:8.2.2]{307} \\
\hyperref[sec:8.2.3]{8.2.3 并交检验与交并检验} & \hyperref[sec:8.2.3]{308} \\
\hyperref[sec:8.3]{8.3 评价检验的方法} & \hyperref[sec:8.3]{310} \\
\hyperref[sec:8.3.1]{8.3.1 错误概率与功效函数} & \hyperref[sec:8.3.1]{310} \\
\hyperref[sec:8.3.2]{8.3.2 最强检验} & \hyperref[sec:8.3.2]{314} \\
\hyperref[sec:8.3.3]{8.3.3 并交检验与交并检验的大小} & \hyperref[sec:8.3.3]{319} \\
\hyperref[sec:8.3.4]{8.3.4 $p$值} & \hyperref[sec:8.3.4]{322} \\
\hyperref[sec:8.3.5]{8.3.5 损失函数意义下的最优性} & \hyperref[sec:8.3.5]{324} \\
\hyperref[sec:8.4]{8.4 习题} & \hyperref[sec:8.4]{326} \\
\hyperref[sec:8.5]{8.5 杂录} & \hyperref[sec:8.5]{335} \\
\hyperref[sec:9]{9 区间估计} & \hyperref[sec:9]{337} \\
\hyperref[sec:9.1]{9.1 引言} & \hyperref[sec:9.1]{337} \\
\hyperref[sec:9.2]{9.2 寻找区间估计量的方法} & \hyperref[sec:9.2]{339} \\
\hyperref[sec:9.2.1]{9.2.1 反演检验统计量} & \hyperref[sec:9.2.1]{339} \\
\hyperref[sec:9.2.2]{9.2.2 枢轴量} & \hyperref[sec:9.2.2]{345} \\
\hyperref[sec:9.2.3]{9.2.3 对累积分布函数作枢轴变换} & \hyperref[sec:9.2.3]{348} \\
\hyperref[sec:9.2.4]{9.2.4 贝叶斯区间} & \hyperref[sec:9.2.4]{352} \\
\hyperref[sec:9.3]{9.3 评价区间估计量的方法} & \hyperref[sec:9.3]{356} \\
\hyperref[sec:9.3.1]{9.3.1 大小与覆盖概率} & \hyperref[sec:9.3.1]{356} \\
\hyperref[sec:9.3.2]{9.3.2 与检验相关的最优性} & \hyperref[sec:9.3.2]{359} \\
\hyperref[sec:9.3.3]{9.3.3 贝叶斯最优性} & \hyperref[sec:9.3.3]{362} \\
\hyperref[sec:9.3.4]{9.3.4 损失函数意义下的最优性} & \hyperref[sec:9.3.4]{363} \\
\hyperref[sec:9.4]{9.4 习题} & \hyperref[sec:9.4]{365} \\
\hyperref[sec:9.5]{9.5 杂录} & \hyperref[sec:9.5]{375} \\
\hyperref[sec:10]{10 渐近评价} & \hyperref[sec:10]{377} \\
\hyperref[sec:10.1]{10.1 点估计} & \hyperref[sec:10.1]{377} \\
\hyperref[sec:10.1.1]{10.1.1 相合性} & \hyperref[sec:10.1.1]{377} \\
\hyperref[sec:10.1.2]{10.1.2 效率} & \hyperref[sec:10.1.2]{379} \\
\hyperref[sec:10.1.3]{10.1.3 计算与比较} & \hyperref[sec:10.1.3]{382} \\
\hyperref[sec:10.1.4]{10.1.4 自助法标准误} & \hyperref[sec:10.1.4]{385} \\
\hyperref[sec:10.2]{10.2 稳健性} & \hyperref[sec:10.2]{388} \\
\hyperref[sec:10.2.1]{10.2.1 均值与中位数} & \hyperref[sec:10.2.1]{389} \\
\hyperref[sec:10.2.2]{10.2.2 M估计量} & \hyperref[sec:10.2.2]{391} \\
\hyperref[sec:10.3]{10.3 假设检验} & \hyperref[sec:10.3]{394} \\
\hyperref[sec:10.3.1]{10.3.1 似然比检验的渐近分布} & \hyperref[sec:10.3.1]{394} \\
\hyperref[sec:10.3.2]{10.3.2 其他大样本检验} & \hyperref[sec:10.3.2]{397} \\
\hyperref[sec:10.4]{10.4 区间估计} & \hyperref[sec:10.4]{401} \\
\hyperref[sec:10.4.1]{10.4.1 近似最大似然区间} & \hyperref[sec:10.4.1]{401} \\
\hyperref[sec:10.4.2]{10.4.2 其他大样本区间} & \hyperref[sec:10.4.2]{403} \\
\hyperref[sec:10.5]{10.5 习题} & \hyperref[sec:10.5]{408} \\
\hyperref[sec:10.6]{10.6 杂录} & \hyperref[sec:10.6]{417} \\
\hyperref[sec:11]{11 方差分析与回归} & \hyperref[sec:11]{421} \\
\hyperref[sec:11.1]{11.1 引言} & \hyperref[sec:11.1]{421} \\
\hyperref[sec:11.2]{11.2 单因素方差分析} & \hyperref[sec:11.2]{421} \\
\hyperref[sec:11.2.1]{11.2.1 模型与分布假设} & \hyperref[sec:11.2.1]{423} \\
\hyperref[sec:11.2.2]{11.2.2 经典的方差分析假设} & \hyperref[sec:11.2.2]{424} \\
\hyperref[sec:11.2.3]{11.2.3 关于均值线性组合的推断} & \hyperref[sec:11.2.3]{426} \\
\hyperref[sec:11.2.4]{11.2.4 方差分析的$F$检验} & \hyperref[sec:11.2.4]{428} \\
\hyperref[sec:11.2.5]{11.2.5 对比的同时估计} & \hyperref[sec:11.2.5]{431} \\
\hyperref[sec:11.2.6]{11.2.6 平方和的分解} & \hyperref[sec:11.2.6]{433} \\
\hyperref[sec:11.3]{11.3 简单线性回归} & \hyperref[sec:11.3]{435} \\
\hyperref[sec:11.3.1]{11.3.1 最小二乘：一种数学解法} & \hyperref[sec:11.3.1]{438} \\
\hyperref[sec:11.3.2]{11.3.2 最优线性无偏估计量：一种统计解法} & \hyperref[sec:11.3.2]{440} \\
\hyperref[sec:11.3.3]{11.3.3 模型与分布假设} & \hyperref[sec:11.3.3]{443} \\
\hyperref[sec:11.3.4]{11.3.4 正态误差下的估计与检验} & \hyperref[sec:11.3.4]{445} \\
\hyperref[sec:11.3.5]{11.3.5 在给定$x=x_0$处的估计与预测} & \hyperref[sec:11.3.5]{450} \\
\hyperref[sec:11.3.6]{11.3.6 同时估计与置信带} & \hyperref[sec:11.3.6]{452} \\
\hyperref[sec:11.4]{11.4 习题} & \hyperref[sec:11.4]{455} \\
\hyperref[sec:11.5]{11.5 杂录} & \hyperref[sec:11.5]{463} \\
\hyperref[sec:12]{12 回归模型} & \hyperref[sec:12]{467} \\
\hyperref[sec:12.1]{12.1 引言} & \hyperref[sec:12.1]{467} \\
\hyperref[sec:12.2]{12.2 变量含误差的回归} & \hyperref[sec:12.2]{467} \\
\hyperref[sec:12.2.1]{12.2.1 函数关系与结构关系} & \hyperref[sec:12.2.1]{469} \\
\hyperref[sec:12.2.2]{12.2.2 一种最小二乘解法} & \hyperref[sec:12.2.2]{470} \\
\hyperref[sec:12.2.3]{12.2.3 最大似然估计} & \hyperref[sec:12.2.3]{472} \\
\hyperref[sec:12.2.4]{12.2.4 置信集} & \hyperref[sec:12.2.4]{476} \\
\hyperref[sec:12.3]{12.3 逻辑回归} & \hyperref[sec:12.3]{478} \\
\hyperref[sec:12.3.1]{12.3.1 模型} & \hyperref[sec:12.3.1]{478} \\
\hyperref[sec:12.3.2]{12.3.2 估计} & \hyperref[sec:12.3.2]{480} \\
\hyperref[sec:12.4]{12.4 稳健回归} & \hyperref[sec:12.4]{483} \\
\hyperref[sec:12.5]{12.5 习题} & \hyperref[sec:12.5]{487} \\
\hyperref[sec:12.6]{12.6 杂录} & \hyperref[sec:12.6]{492} \\
\hyperref[app:computer-algebra]{计算机代数} & \hyperref[app:computer-algebra]{497} \\
\hyperref[app:distributions]{常见分布表} & \hyperref[app:distributions]{503} \\
\hyperref[app:references]{参考文献} & \hyperref[app:references]{509} \\
\hyperref[app:index]{索引} & \hyperref[app:index]{525} \\
\end{longtable}\endgroup
